% !TEX root = main.tex
%
% Cross-volume convention firewall.
%
% This appendix is only a boundary condition for reading the local
% theorem.  It is not a transfer dictionary and it records no export
% profile into neighbouring programmes.

\section{Cross-volume convention firewall}\label{app:convention-contract}

The theorem proved in this manuscript is local on
\[
  \R^2_{\mathrm{top}}\times\widehat{\C^2}_{0,\mathrm{hol}}
  =
  \R_{\mathrm{brane}}\times\R_{\mathrm{top,normal}}
  \times\widehat{\C^2}_{0,\mathrm{hol}},
\]
with brane line \(\R_{\mathrm{brane}}\times\{0\}\times\{0\}\). Its
closed sector is Hamiltonian BF theory on the formal
holomorphic-symplectic disk together with the topological directions, and
its open sector is the Dirac brane probe on the brane line. The statement
is at \(d=\dim_\C Y=2\) in the holomorphic factor and uses the brane line
\(\R_{\mathrm{brane}}\) as its one-dimensional factorization direction.

This appendix fixes the boundary of that statement.  It does not
define a cross-volume dictionary.  It does not record forward or
backward transfer profiles.  It only states what the local theorem
does not assert.

\begin{rmk}[Non-assertions]
  No compact Calabi--Yau threefold theorem is asserted here.  No
  theorem of Vol III is asserted here.  No Igusa cusp-form, Borcherds
  product, or BKM denominator identity is asserted here.  No global
  twisted-M-theory realization is asserted here.  No global
  CY-to-chiral statement, compact-CY$_3$ chiral-homology statement, or
  global BCOV amplitude statement follows from the local Hamiltonian
  BF/Moyal calculation.
\end{rmk}

\begin{rmk}[Allowed use]
  The local theorem may be used as a convention check: it fixes the
  $d=2$ formal-disk normalization, the brane-line factorization
  convention, the scalar anomaly class $[\bar c]$, and the reduced
  degree-zero Moyal/Capelli coefficients.  This is not a compact
  CY$_3$ trace pairing, not an $S^3$ framing, not a specialisation to
  a chiral curve, and not an automorphic or Borcherds construction.
\end{rmk}

\begin{rmk}[Matched-conventions requirement]
  Any global transfer to Vol III, Vol II, Vol I, the Igusa cusp-form
  programme, a compact Calabi--Yau threefold, a $K3\times E$
  specialisation, a compact-CY$_3$ chiral-homology theorem, a global
  BCOV amplitude theorem, or a BKM denominator identity requires a
  separate matched-conventions theorem.  Such a theorem would have to
  state its holomorphic dimension, worldsheet or specialisation curve,
  trace pairing, framing datum, coefficient topology, anomaly
  normalization, deformation parameter, QME hypothesis, exact comparison
  morphism, and claim status.  No such theorem is proved or assumed in
  this manuscript, and no sibling-volume theorem may use the local result
  or this appendix as a substitute for these data.
\end{rmk}

\begin{lem}[No formal-disk transfer]\label{lem:no-formal-disk-transfer}
  The local theorem at
  \(\R^2_{\mathrm{top}}\times\widehat{\C^2}_{0,\mathrm{hol}}\), with brane
  line \(\R_{\mathrm{brane}}\times\{0\}\times\{0\}\), does not determine a
  compact Calabi--Yau threefold theorem, a global BCOV theorem, a Vol
  III CY-to-chiral or chiral-homology theorem, an Igusa cusp-form
  theorem, a Borcherds product theorem, a BKM denominator theorem, or a
  \(K3\times E\) specialization.
\end{lem}

\begin{proof}
  The local input consists of the formal Hamiltonian algebra
  \(\C[[z_1,z_2]]/\C\), its continuous cotangent extension, the point
  brane Ext algebra, the brane-line factorization-current convention,
  the scalar quotient, and the Capelli/Weyl normalization.  These data
  are invariant under the formal Darboux germ and omit compact topology,
  a negative-cyclic Calabi--Yau class, a worldsheet or specialization
  curve, an \(S^3\) framing, a global BV propagator, the holomorphic
  anomaly equation, Borcherds input data, a Weyl chamber, a compact Hall
  orientation, and a compact chiral-homology object.  Hence the local
  theorem can serve only as a \(d=2\) convention check until a separate
  matched-conventions theorem supplies the missing global data.
\end{proof}

\begin{defn}[Matched-conventions theorem datum]\label{defn:matched-conventions-theorem-datum}
  A cross-volume transfer theorem citing the local result must state at
  least the following datum: the holomorphic dimension, the ambient
  geometric object, the worldsheet or factorization-homology
  specialization curve, the brane or boundary condition, the trace
  pairing, the framing, the coefficient topology, the anomaly
  normalization, the deformation parameter, the QME or obstruction
  hypothesis, the exact comparison morphism, the claim status, and the
  precise role of the local theorem as convention check, source theorem,
  or unused comparison datum.
\end{defn}

\begin{rmk}[Templates, not theorems]\label{rmk:matched-conventions-templates}
  The following are templates for future separate theorems.
  They are not assertions in this manuscript.
  \begin{enumerate}[label=(\roman*)]
    \item \emph{Vol III CY-to-chiral template.}  State \(d\), the
      compact or non-compact CY object, the CY class in negative cyclic
      homology, the stage-one object \(\Phi^{\mathrm{FA}}_d\), the
      stage-two specialization \((\Sigma_{d-1},C)\), the trace pairing,
      the framing, and the QME/anomaly witness.  Then prove that the
      \(d=2\) formal-disk normalization used here is the stated local
      convention check for that theorem.
    \item \emph{Compact BCOV template.}  State the compact CY$_3\),
      holomorphic volume form, BCOV field complex, propagator,
      gauge-fixing, boundary or brane condition, holomorphic anomaly
      equation, zero-mode treatment, and counterterm scheme.  Then prove
      the global amplitude statement inside that BCOV theory.  The
      local Moyal coefficient can serve only as a boundary normalization
      after these choices are fixed.
    \item \emph{Igusa/BKM template.}  State the Jacobi or modular input,
      Borcherds lift normalization, lattice and Weyl chamber, character,
      compact \(K3\times E\) or CHL host if one is claimed, Hall or
      chiral source, orientation line, primitive recognition theorem,
      and scalar DT/PT/OP normalization.  Then prove the denominator or
      compact realization theorem in that setting.  The local theorem
      supplies no Borcherds input form, no Weyl denominator, and no
      compact Hall orientation.
  \end{enumerate}
  These templates are failure tests.  If the listed datum is absent, the
  proposed cross-volume sentence is not a theorem and must be read as an
  open obligation.
\end{rmk}

% --- End of appendix.
